% wg-findings.tex v1.0 2026-08-31 — HAnim WG contribution packet: AllBonesLOA5Skeletons findings + differential-probe rediscovery (Web3D 2026 Doha)
\documentclass[10pt]{article}
\usepackage[margin=1.9cm]{geometry}
\usepackage[T1]{fontenc}
\usepackage{lmodern}
\usepackage{microtype}
\usepackage{booktabs}
\usepackage{enumitem}
\usepackage[hidelinks]{hyperref}
\usepackage{xcolor}
\usepackage{titlesec}

\titleformat{\section}{\normalfont\large\bfseries}{\thesection.}{0.6em}{}
\titlespacing*{\section}{0pt}{1.1ex plus 0.4ex}{0.6ex}
\setlist{nosep,leftmargin=1.4em}
\newcommand{\code}[1]{\texttt{#1}}
\setlength{\parskip}{0.35ex}
\setlength{\parindent}{1.2em}

\title{\vspace{-1.4em}\textbf{Findings in \code{AllBonesLOA5Skeletons}, with a mechanical rediscovery}\vspace{-0.4em}}
\author{Alexander H. Hoffman\\\normalsize \href{mailto:alex@euphe.me}{alex@euphe.me} \,\textperiodcentered\, EUPHEME Technologies LLC}
\date{\normalsize Prepared for the HAnim Working Group, Web3D 2026, Doha \,\textperiodcentered\, August 2026}

\begin{document}
\maketitle
\vspace{-1.6em}

\noindent\textbf{Spirit of this note.} The \code{AllBonesLOA5Skeletons} set is a gift: 277
mesh files in which every one of the 256 leaf parts already carries a \code{TouchSensor},
a DEF'd material, a hidden close-up \code{Viewpoint}, and hand-written anatomical prose in
\code{<meta name='description'>} (275 of 277). The interactivity for a full anatomy explorer
was authored into these assets years ago. While building such an explorer --- making every
bone individually addressable --- we surfaced a small number of file-level defects and two
tooling traps. Everything below is reproducible from the published meshes; concrete fixes
are proposed for each item. \textit{[TODO: insert upstream GitHub issue numbers once filed.]}

\section{\code{AxesDisplay.x3d}: referenced by 260 meshes, shipped with none}
260 of the 277 mesh files contain
\code{<Inline DEF='AxesDisplay' url='"AxesDisplay.x3d"'/>}, and the file does not exist in
the set. Unlike every other \code{Inline} in these files, this \code{url} carries
\emph{no fallback list} --- the others all pair a local path with a
\code{https://www.web3d.org/...} mirror. Loading the assembled skeleton therefore fires
\textbf{260 hard 404s} before the scene is usable.

\emph{Proposed fix:} ship \code{AxesDisplay.x3d} alongside the meshes, or give this
\code{Inline} the same web fallback URL the other inlines already have.

\section{\code{l\_tarsal\_distal\_phalanx\_5.x3d} contains no geometry}
The file exists (7{,}630 bytes) but holds zero \code{IndexedFaceSet} and zero
\code{Coordinate} nodes --- header, \code{TouchSensor}, and \code{Transform} boilerplate
around nothing. Its right-side twin is a complete mesh:

\begin{center}\small
\begin{tabular}{@{}lrrr@{}}
\toprule
file & bytes & \code{IndexedFaceSet} & \code{Coordinate}\\
\midrule
\code{l\_tarsal\_distal\_phalanx\_5.x3d} & 7{,}630 & 0 & 0\\
\code{r\_tarsal\_distal\_phalanx\_5.x3d} & 11{,}001 & 2 & 1\\
\bottomrule
\end{tabular}
\end{center}

The left foot renders with 13 phalanges to the right foot's 14; the assembled skeleton has
198 bones where it should have 199, and the left little toe visibly lacks its tip. (A
bilateral fusion of the fifth toe's middle and distal phalanges is a real anatomical
variant --- but it is bilateral. A one-sided absence is an asset defect, not a variant.)

\emph{Proposed fix:} regenerate the mesh, or mirror \code{r\_tarsal\_distal\_phalanx\_5.x3d}.

\section{Not a bug: the ethmoid --- a trap for tooling, recorded because we fell in}
The set contains \code{ethmoid.x3d}, \code{l\_ethmoid.x3d}, and \code{r\_ethmoid.x3d}. The
ethmoid is one unpaired midline bone, so three files look like a double-count --- and any
tool that decides \emph{what a bone is} from filenames will count a cranium of 9 against
the canonical 8. The assets are in fact \textbf{correct}: \code{loa5\_humanoid.x3dfrag}
inlines \code{ethmoid.x3d} exactly once, and that file is a geometry-free container
inlining the two lateral labyrinths. The trap is that this container is structurally
identical to the LOA-1 \emph{region} containers (\code{skull.x3d} inlines 22 bones;
\code{l\_carpal.x3d} inlines 4) yet means something different: a region's children are
separate bones, the ethmoid's children are pieces of one bone. A usable discriminator: a
container is a split \emph{bone} when every child's name reduces to the container's own
name once the \code{l\_}/\code{r\_} prefix is stripped. Across all 277 meshes the ethmoid
is the only such case.

\emph{No fix needed;} a short comment inside \code{ethmoid.x3d} saying so would spare the
next tool author.

\section{The auditory ossicles and the hyoid are absent}
No \code{malleus}, \code{incus}, \code{stapes}, or \code{hyoid} mesh exists under any name:
the string ``hyoid'' appears nowhere in the 277 mesh files nor in
\code{loa5\_humanoid.x3dfrag}. That is 7 of the 206 bones of the adult skeleton. The
ossicles are arguably out of scope for a surface skeleton --- they live inside the temporal
bone. The \textbf{hyoid is not}: it is a visible, palpable, freestanding bone of the neck,
anatomically notable precisely because it articulates with no other bone, and its absence
is conspicuous in any scene showing the mandible and cervical spine.

\emph{Proposed fix:} add \code{hyoid.x3d}; consider documenting the ossicles as a
deliberate omission.

\section{A processing trap: rewriting \code{<Material USE=.../>} whitens the teeth}
Not an asset bug, but a hazard for anyone modernizing the set. A regex such as
\begin{center}
\code{re.sub(r"<Material\textbackslash b[\textasciicircum>]*/>", ...)}
\end{center}
also matches
\code{<Material USE='ToothMaterial'/>}; rewriting a \code{USE} reference as a definition
severs it from its \code{DEF}, and all 32 teeth (which share one \code{ToothMaterial} via
\code{USE}) lose their material and render pure white. 34 meshes still carry legacy Phong
\code{<Material>} (the 32 \code{tooth\_*} plus \code{l\_/r\_ethmoid}); converting them to
\code{PhysicalMaterial} is worthwhile --- they blow out under PBR studio lighting --- but
the conversion must retag \code{USE} references without touching their attributes, never
rewrite them as definitions.

\section{The zero-geometry defect class is mechanically detectable}
Finding 2 was first noticed by eye. It need not be. A \emph{differential object-presence
probe} --- render the scene with a node, render it with that node removed, and count the
changed pixels --- attributes to every node its per-node pixel contribution. Run against a
real renderer (X\_ITE 15.1.12 under headless Chromium, 512$\times$384 frames, recorded run
of 2026-08-01), a correct \code{HAnimHumanoid} contributes \textbf{7{,}744 px} to the frame;
a defective one contributes \textbf{0 px} --- inside a frame that is \emph{not} blank
(the pedestal still contributes 7{,}068 px, so a whole-frame blankness gate passes it).
The probe's control run is clean: rendering an identical scene twice differs by 0 px, and
the contribution threshold (98 px) sits far above that noise floor.

This is exactly the signature of the empty phalanx: a node present in the document,
absent from the image, in a scene that otherwise renders. The same probe therefore
rediscovers Finding 2 mechanically, with no anatomical knowledge required:
\code{l\_tarsal\_distal\_phalanx\_5} shows a 0 px contribution while its right-side
twin contributes normally. The practical point for the working group: \textbf{this defect class
is CI-detectable}. A per-mesh presence probe over the asset set (each mesh rendered with
and without its geometry payload) would have flagged the empty file at publication time,
and would guard every future regeneration of the set. We are glad to contribute the probe
and the runbook.

\vspace{0.5ex}
\noindent\textbf{Closing.} None of this diminishes the set --- the craft was already
authored in; these are finishing touches. We offer the findings, the proposed fixes, and
the probe in the same spirit the assets were offered: as a contribution to the commons.
\textit{[TODO: GitHub issue numbers here once filed.]}

\end{document}
